% compendium/execution-market-making/market-organisation.tex
% Entries migrated on 2026-10-16 from the first five files of the legacy foundations folder (slice 35), owned by
% the section overview execution-market-making, one of the files' two map candidates. Dispositions in
% notes/migration-slices.md; the files verbatim in notes/artefacts/execution-market-making/.

\section{Markets, traders and venues}

\begin{qfdefinition}[What markets do: liquidity, price discovery, risk transfer and capital allocation; private and public benefits; why the mechanism matters]\label{def:market-functions}
\qffield{Statement} Harris organises the study of trading around five characteristics of market quality: liquidity, transaction costs, informative prices, volatility and trading profits, and the secondary question of how market structure, ``trading rules and information systems'', affects each \citep[chap.~1, sec.~1.2, p.~4]{alpha-harris}. What markets produce falls into private benefits, which accrue to traders when they trade (they trade voluntarily, so they expect to gain), and public benefits, positive externalities that reach people who do not trade, in two classes: those from informative prices and those from liquid markets \citep[chap.~9, secs.~9.1.1--9.3, pp.~205--206]{alpha-harris}. The functions the legacy notes listed map onto Harris's account as follows:
\begin{itemize}
\item \emph{liquidity}: the buy side buys exchange services, of which liquidity, ``the ability to trade when you want to trade'', is the most important, and the sell side sells it \citep[chap.~3, sec.~3.1, p.~32]{alpha-harris}; patient traders get better prices, and impatient traders ``pay for the privilege of trading when they want to trade'' (chap.~1, p.~5); the dimensions are in \hyperref[def:liquidity-dimensions]{liquidity and market quality};
\item \emph{price discovery}: markets ``aggregate information from many different sources to produce prices that are closer to fundamental values than anyone could consistently estimate by himself or herself'' (chap.~9, sec.~9.4, p.~206); how informed traders do this is \hyperref[res:informed-trading-aggregation]{informed trading and aggregation};
\item \emph{risk transfer}: hedgers trade to exchange risks (chap.~8, sec.~8.1.3, pp.~182--188); liquid hedging and insurance markets let producers specialise, and risk sharing through public corporations lets large projects be financed, which lowers costs of capital and prices even for those who never trade (chap.~9, secs.~9.5.2--9.5.3, pp.~214--215);
\item \emph{capital allocation}: informative primary-market prices decide which new projects receive capital, and informative secondary-market prices help investors choose, pay and replace managers of existing capital (chap.~9, secs.~9.4.2--9.4.3, pp.~211--212); when prices do not reflect how well managers perform, good managers may be dismissed and poor ones kept (p.~211).
\end{itemize}
\emph{The mechanism matters.} Hasbrouck: ``The market-clearing price, at least at it arises in usual Walrasian tatonnement, rarely appears in microstructure analyses''; at one instant there are many prices, by direction, urgency, identity and relationship to the counterparty \citep[chap.~1, sec.~1.1.3, p.~4]{s6-hasbrouck}. Harris: market structure ``determines what people can know and do in a market'' and so the power relations among traders \citep[chap.~5, p.~89]{alpha-harris}.

\qffield{Conditions} The list of functions is a classification, not a theorem; Harris adds that markets work well only when negotiation costs are small relative to the value traded, so economies mix markets with firms, which are small command economies (chap.~9, p.~209). Which benefit dominates is an empirical matter for each market. The welfare ranking of these functions is Harris's opinion, which he labels as such: public policy should first serve utilitarian traders, then the public benefits of liquid, informative markets, and support profit-motivated traders only as far as they serve those (sec.~9.6, pp.~216--217); see \hyperref[def:market-regulation-design]{regulation as market design}.

\qffield{Practice} \emph{Trade-offs.} Objectives conflict: rules that restrict insider trading ``usually increase liquidity while making prices less informative'' (chap.~9, p.~218); and the transparency that improves price formation is what each trader wants to avoid revealing (Lehalle and Laruelle, quoted in \hyperref[def:market-structures-venues]{market structures and venues}). \emph{Visible activity against economic importance.} Harris's own estimates for the United States (Table~3-4, ``Author's estimates'', c.\ 2000) put common stocks at about 20 percent of capital wealth and real estate at about 50 percent; derivatives represent no wealth because they are in zero net supply; ``Stocks represent less wealth than the widespread attention given to them by the media would suggest'' \citep[chap.~3, sec.~3.4.1, p.~45]{alpha-harris}. The figures are dated rough estimates; the point is that trading volume and attention are not measures of a market's economic weight.

\qffield{Pitfalls} The legacy principle ``immediacy has a price'' said the spread compensates traders who absorb ``inventory, funding, processing, and adverse-selection risks''; the components and their identification are \hyperref[res:spread-determinants]{spread determinants} and \hyperref[res:spread-component-identification]{spread components}, and ``funding'' is the legacy notes' addition, not in Harris's list as checked there. The legacy statement that liquid markets ``lower financing costs'' is Harris's argument through investors' required returns (p.~215), not a measured effect. Harris's claim that liquidity benefits even non-traders through cheaper goods is an argument about externalities, stated without evidence in these pages.

\qffield{See also} \hyperref[res:zero-sum-skill-luck]{zero-sum trading, skill and luck}; \hyperref[def:trader-motives]{trader motives}; \hyperref[def:liquidity-dimensions]{liquidity and market quality}; \hyperref[def:market-efficiency]{market efficiency}; \hyperref[def:market-regulation-design]{regulation as market design}.

\qffield{Further reading} Harris, chaps.~1 and~9; Hasbrouck, chap.~1. The legacy file is kept verbatim in \nolinkurl{notes/artefacts/execution-market-making/}.

\qffield{Index} market functions; market quality; liquidity provision; price discovery; risk transfer; risk sharing; capital allocation; primary markets; secondary markets; private benefits of trading; public benefits; positive externalities; Walrasian price; market structure; stock market wealth share

\qffield{Status} Verification: partly checked 2026-10-16 (against the text layers of local copies of \texttt{alpha-harris}, 2002 Oxford printing, pp.~3--9, 32, 45, 89 and 202--218, and \texttt{s6-hasbrouck}, 2007 Oxford printing, pp.~3--6). Stability: draft. Currency: concepts timeless; Harris's wealth shares are estimates for about 2000.
\end{qfdefinition}

\begin{qfdefinition}[Trader motives: utilitarian, profit-motivated and futile traders; patience and the liquidity side; instruments for hedging; sources of a trading edge]\label{def:trader-motives}
\qffield{Statement} Harris classifies traders by why they trade, using stylised traders who each trade for one reason \citep[chap.~8, p.~177]{alpha-harris}:
\begin{itemize}
\item \emph{utilitarian traders} trade for a benefit other than trading profit: investors and borrowers move money forward or backward through time, asset exchangers trade one asset for another of greater immediate use, hedgers trade risks, gamblers trade for entertainment, and also fledglings (to learn), cross-subsidisers (to pay for broker services through commissions) and tax avoiders (sec.~8.1, p.~178; Table~3-1, p.~33);
\item \emph{profit-motivated traders} trade only because they rationally expect to profit: speculators, who are informed traders or parasitic traders (order anticipators and bluffers), and dealers, who sell liquidity (sec.~8.2, pp.~194--196, Table~8-4);
\item \emph{futile traders} ``believe that they are profit-motivated traders. Although they expect to trade profitably, their expectations are not rational'' (p.~177): inefficient traders, of whom the commonest are pseudo-informed traders trading on old news, victimised traders whose agents fail their fiduciary duties, and rogue traders who hide losses from employers or clients (sec.~8.3, pp.~197--198).
\end{itemize}
A trader with a \emph{long} position owns the instrument and profits when prices rise; one with a \emph{short} position ``ha[s] sold something that they do not own'' and repurchases, or \emph{covers}, later (chap.~3, p.~32). The legacy's ``or has equivalent negative exposure'' widens Harris's definition to derivative positions.

\qffield{Conditions} The legacy ``trading game'' view, players submitting orders under venue rules, is Harris's point that market structure ``determines what people can know and do'' and so ``the power relationships among different types of traders'' (chap.~1, p.~7; chap.~5, p.~89). Motives are a modelling device: ``In practice, traders often trade for many reasons'', and Harris's sidebar \emph{Multiple Identities} says many traders ``simultaneously invest, speculate, and gamble'' (p.~177). Motive is therefore inferred from context and exposure, not from the order: Harris notes that ``Order type therefore is not always the best indicator of whether a trader is offering liquidity or taking it'' \citep[chap.~19, sec.~19.3.1, p.~400]{alpha-harris}. \emph{Patience and the liquidity side}: impatient traders search actively and demand liquidity; patient traders search passively and supply it (sec.~19.1.2, p.~396); Hasbrouck adds that with electronic access the role ``is a strategic choice that can be quickly reversed'' \citep[chap.~1, sec.~1.2, p.~5]{s6-hasbrouck}.

\qffield{Practice} \emph{Instruments for a problem.} Investors favour high expected return, low risk and low transaction cost; borrowers the cheapest instruments to create, service and repay (chap.~8, p.~180). Hedgers like liquid markets and instruments closely correlated with their risks (sec.~8.1.3.4, p.~188). Harris's hedging examples give the vehicle trade-offs: forwards need the parties to find and trust each other and agree delivery terms, so they suit traders who deal with each other in the normal course of business and ``forward markets generally are not liquid'' (pp.~184--185); futures are standardised forwards guaranteed by a clearinghouse, liquid, and a \emph{linear} hedge that removes upside as well as downside (p.~185), with the delivery location differing from the hedger's (pp.~185--186); options are a \emph{nonlinear} hedge that limits downside for a premium and keeps upside (pp.~186--187); swaps exchange fixed for floating cash flows and leave basis risk, since a bank's deposit rates are not the swap's index rate (pp.~187--188). Basis risk is developed in \hyperref[def:futures-hedge-basis]{futures hedging and basis}. \emph{Comparative advantage.} ``Successful profit-motivated traders must have some advantage that allows them to trade profitably. Their various advantages usually involve information'' (p.~196), and free entry drives the net profits of alternative strategies towards equality after all costs (chap.~1, pp.~7--8). Harris advises speculators to expose themselves ``only to the risks they understand best'' (p.~186). The legacy list of edges (information, model, execution, risk bearing, discipline) is a practitioner synthesis; only information and the cost of trading are Harris's in these pages.

\qffield{Pitfalls} The legacy taxonomy put arbitrageurs among informed traders who ``supply or demand liquidity''; Harris counts arbitrageurs as informed traders about relative values (Table~8-4) and, as liquidity suppliers, calls them ``porters of liquidity'' (chap.~19, p.~404); see \hyperref[def:arbitrage]{arbitrage}. Noise traders, in the legacy definition ``rebalancing, cash needs, index changes, behavioural demand, or errors'', mix utilitarian traders with futile ones; Harris's point is that utilitarian and futile traders lose on average to profit-motivated traders ``because trading is a zero-sum game'' (p.~177). The legacy futile list put ``pseudo-informed'' beside ``inefficient''; in Harris the first is a kind of the second. Gamblers are utilitarian in Harris's scheme because they get entertainment, not because of ``payoff convexity'' as the legacy notes said.

\qffield{See also} \hyperref[res:zero-sum-skill-luck]{zero-sum trading, skill and luck}; \hyperref[def:informed-trader-types]{informed trader types}; \hyperref[def:order-anticipation]{order anticipation}; \hyperref[def:bluffing-manipulation]{bluffing and manipulation}; \hyperref[def:market-intermediaries]{market intermediaries}; \hyperref[def:market-functions]{what markets do}.

\qffield{Further reading} Harris's chapter~8 is the source of the classification. The legacy blocks it replaces are kept verbatim, with the other blocks of their file, in the folder \nolinkurl{notes/artefacts/execution-market-making/}.

\qffield{Index} trader motives; utilitarian traders; profit-motivated traders; futile traders; investors; borrowers; asset exchangers; hedgers; gamblers; fledglings; cross-subsidisers; speculators; parasitic traders; pseudo-informed traders; victimised traders; rogue traders; long position; short position; covering; patience; linear hedge; nonlinear hedge; forwards; futures; options; swaps; comparative advantage; noise traders

\qffield{Status} Verification: partly checked 2026-10-16 (against the text layers of local copies of \texttt{alpha-harris}, 2002 Oxford printing, pp.~5--8, 32--33, 176--201, 396 and 400--404, and \texttt{s6-hasbrouck}, 2007 Oxford printing, p.~5). Stability: draft. Currency: timeless (classification).
\end{qfdefinition}

\begin{qfresult}[Zero-sum trading accounting, its two corollaries, and what past profit says about skill]\label{res:zero-sum-skill-luck}
\qffield{Statement} \emph{Accounting.} Harris: trading is ``a zero-sum game in an important accounting sense'': ``the combined gains and losses of buyers and sellers always sum to zero. If a buyer profits from a trade, the seller loses the opportunity to profit by the same amount'' \citep[chap.~8, p.~176]{alpha-harris}; elsewhere, ``when gains and losses are measured relative to the market average'' (chap.~1, p.~6). For a trade of $q$ at price $P$ and any common terminal value $V$, the buyer's gain $q(V-P)$ and the seller's $q(P-V)$ sum to zero; with costs $C_i\ge0$ the sum of net gains is $-\sum_iC_i$. \emph{Corollary 1}: ``Since trading is a zero-sum game, speculators and dealers cannot profit on average if they trade only among themselves. They can profit only if utilitarian traders are willing to trade'' (chap.~9, p.~206). \emph{Corollary 2} (the arithmetic of active management, derived here, not from Harris): if active managers' holdings aggregate, value-weighted, to the market portfolio, their aggregate gross active return is zero and their aggregate net active return is $-\sum_iw_ic_i\le0$; a single manager beats a passive alternative with drag $c_P$ only if $a_i>c_i-c_P$. \emph{Skill and luck} (restated in this slice): write profit as $\Pi_{t+1}=S_{t+1}+L_{t+1}$, skill plus luck, and let $\mathcal F_t$ be the evaluator's information. Then
\[ E[\Pi_{t+1}\mid\mathcal F_t,\Pi_t]=E[\Pi_{t+1}\mid\mathcal F_t]\quad\text{if}\quad E[S_{t+1}\mid\mathcal F_t,\Pi_t]=E[S_{t+1}\mid\mathcal F_t]\ \text{and}\ E[L_{t+1}\mid\mathcal F_t,\Pi_t]=E[L_{t+1}\mid\mathcal F_t]. \]
Both conditional conditions are needed. Unconditional zero-mean luck, $E[L]=0$, is not enough: persistent luck makes past profit predictive without any skill.

\qffield{Conditions} \emph{Accounting.} The identity needs a common valuation of the asset for both sides (a common mark or terminal value); it says nothing about utility, so a hedger who pays to shed risk or an asset exchanger can gain in welfare while losing in money (chap.~8, pp.~176--177). Corollary~1 needs a closed set of traders and costs that are not income of the group (fees paid to members are transfers within it). Corollary~2 needs the actives to hold the market in aggregate, which fails if other holders (for example, non-market holdings by passive owners or issuers) take the other side. \emph{Skill and luck.} The legacy proposition assumed $S_{t+1}$ independent of past profit given current information and ``zero-mean luck''; its proof used $E[L_{t+1}\mid\mathcal F_t,\Pi_t]=E[L_{t+1}\mid\mathcal F_t]$ without stating it. If skill is \emph{not} in $\mathcal F_t$, past profit is informative about it: with persistent skill $S\sim N(\mu,\tau^2)$ and i.i.d.\ luck of variance $\sigma^2$, the regression slope of $\Pi_{t+1}$ on $\Pi_t$ is $\tau^2/(\tau^2+\sigma^2)$, which is the shrinkage a performance evaluator applies. Harris states the evaluator's version: predicting from past performance needs past performance to reflect skill, the skill to keep working in future market conditions, and the manager to keep it \citep[chap.~22, sec.~22.3, p.~451]{alpha-harris}; a statistical test works because luck ``is random and completely unpredictable'' and its net effect on average returns shrinks with time, while skill has a systematic effect (sec.~22.4.1, p.~454).

\qffield{Practice} Checked (\nolinkurl{slice35-checks.py}, numpy): 200{,}000 simulated managers with $\tau=0.5$, $\sigma=1$ give a slope of 0.200 against $\tau^2/(\tau^2+\sigma^2)=0.200$; with no skill and i.i.d.\ luck the slope is $-0.006$; with no skill and AR(1) luck of persistence 0.5 and unconditional mean zero the slope is 0.501, so past profit predicts luck; with skill in the information set and conditionally mean-zero luck the incremental slope is 0.0005. Random bilateral trades among ten traders marked at a common value sum to zero gross ($-7\times10^{-13}$) and to minus total costs net; value-weighted active returns of managers aggregating to the market sum to zero, and do not when they do not. In practice: test skill claims against a benchmark, over long samples (Harris notes more than 20 years of returns are typically needed for a reliable test, p.~453), and correct for survivorship (sec.~22.6, pp.~470--471).

\qffield{Pitfalls} ``Zero-sum'' is an accounting statement about money, not a claim that trade creates no value. Harris's two wordings differ: relative to a common terminal value (p.~176) and ``relative to the market average'' (p.~6); the second makes it a statement about excess returns. The legacy caveat that a skill claim must explain ``why the decision was right before the outcome was known, why the edge should persist, and why it is not already competed away'' matches Harris's three conditions and his free-entry argument (chap.~1, pp.~7--8). A proposition that conditions on skill being known cannot be used to say that past profit is useless for an evaluator who does not know it.

\qffield{See also} \hyperref[def:trader-motives]{trader motives}; \hyperref[meth:relative-performance-evaluation]{relative performance evaluation}; \hyperref[meth:backtest-evaluation-diagnostics]{backtest evaluation diagnostics}; \hyperref[def:market-functions]{what markets do}.

\qffield{Further reading} Harris, chaps.~8, 9 and~22. The legacy blocks are kept verbatim in \nolinkurl{notes/artefacts/execution-market-making/}; the checks in \nolinkurl{slice35-checks.py} in the same folder.

\qffield{Index} zero-sum game; trading profits; utilitarian traders; profit-motivated traders; arithmetic of active management; active return; fee drag; skill; luck; performance persistence; shrinkage; conditional mean independence; persistent luck; survivorship bias

\qffield{Status} Verification: partly checked 2026-10-16 (against the text layer of a local copy of \texttt{alpha-harris}, 2002 Oxford printing, pp.~4--8, 176--177, 206, 444--454 and 470--471; the corollaries and the skill--luck conditions derived and simulated). Stability: draft. Currency: timeless.
\end{qfresult}

\begin{qfdefinition}[Market intermediaries: dealers and market makers, brokers and their services, broker-dealers, and the suppliers of liquidity]\label{def:market-intermediaries}
\qffield{Statement} \emph{Dealers} trade with their clients when the clients want to trade and profit by buying low and selling high; \emph{brokers} arrange trades as agents and profit from commissions; firms that do both are \emph{broker-dealers} or dual traders \citep[chap.~3, sec.~3.1.2, pp.~33--34]{alpha-harris}. ``The sell side exists only because the buy side will pay for its services'' (p.~34). Dealers sell liquidity at their bid and ask; to profit they must buy and sell equal volumes and so must find the prices that balance supply and demand; Harris classes them as uninformed about fundamental values. Market makers provide liquidity on demand in small quantities and trade in and out many times a day; block facilitators provide liquidity to large traders, may take days or weeks to unwind, and generally choose their clients \citep[chap.~8, sec.~8.2.2, pp.~195--196]{alpha-harris}. \emph{Why brokers are cheaper than trading directly}: they solve clearing and settlement problems, give access to exchanges and dealers their clients cannot reach, know better who might trade, negotiate better, and represent orders when clients are absent \citep[chap.~7, sec.~7.1, pp.~140--141]{alpha-harris}; ``The most important, but perhaps least appreciated'' of these is clearing and settlement, because order-driven markets arrange trades among strangers (p.~141).

\qffield{Conditions} \emph{Who supplies liquidity.} Harris's five suppliers, by what they know and the dimension they supply \citep[chap.~19, secs.~19.3.2--19.3.7, pp.~401--405]{alpha-harris}: market makers supply immediacy to small traders and avoid informed traders; block dealers supply depth to clients they know to be uninformed; value traders, ``the ultimate suppliers of liquidity'', supply depth and make markets resilient, and lose only to better news or value traders (p.~403); precommitted traders post aggressive limit orders to lower the cost of trades they intend anyway and supply immediacy at narrow spreads, which can drive dealers out (p.~404); arbitrageurs are ``not suppliers of liquidity but porters of liquidity'' between markets (p.~404). Harris's short list: ``Impatient traders take liquidity. Dealers, limit order traders, and some speculators offer liquidity. Brokers and exchanges organize liquidity'' (p.~394). Capital standards for dealers only ensure that they \emph{can} supply liquidity; ``no amount of capital will make them willing'' under heavy adverse selection (p.~402).

\qffield{Practice} \emph{Brokered transactions.} Harris's Table~7-1 distinguishes order-flow, block, new and seasoned offering, and merger-and-acquisition business (p.~140). Order-driven exchanges ``are essentially brokerages'' (chap.~3, p.~35), and automated order-driven systems are ``essentially electronic brokers'' (chap.~7, p.~139); regulators classify similar systems sometimes as exchanges and sometimes as brokers (p.~139). Introducing brokers pass orders to clearing members, and sponsored subscribers of electronic systems trade under their broker's real-time credit limits (pp.~142--143). Inventory risk and quote management are \hyperref[def:dealer-inventory-risk]{dealer inventory risk}; agency problems are \hyperref[def:broker-agency-conflicts]{broker agency conflicts}; block roles are \hyperref[def:block-trade]{block trades}; fast intermediaries are \hyperref[def:hft-strategies]{high-frequency trading strategies}.

\qffield{Pitfalls} The legacy taxonomy of liquidity suppliers is Harris's five types with the losses he names; its warning that calling all passive flow ``market making'' hides the source of depth paraphrases his point that each supplier occupies a niche (p.~405). The legacy broker service list's post-trade item adds confirmation, custody, margin and records to Harris's clearing and settlement, and its ``information control'' item is his block-trading material (\hyperref[def:block-trade]{block trades}); in Harris custody is a service of depositories and custodians (chap.~3, sec.~3.2.2.4, p.~38). Dealer compensation from ``rebates'' and ``information'' is the legacy notes' addition.

\qffield{See also} \hyperref[def:trader-motives]{trader motives}; \hyperref[def:market-structures-venues]{market structures and venues}; \hyperref[def:liquidity-dimensions]{liquidity and market quality}; \hyperref[def:outside-spread]{the outside spread}; \hyperref[mod:inventory-based-quotes]{inventory-based quotes}.

\qffield{Further reading} Harris, chaps.~7, 13 and~19. The legacy file is kept verbatim in \nolinkurl{notes/artefacts/execution-market-making/}.

\qffield{Index} dealers; market makers; block facilitators; brokers; broker-dealers; dual traders; sell side; buy side; brokered transactions; introducing brokers; clearing members; sponsored access; liquidity suppliers; precommitted traders; value traders; porters of liquidity; dealer capital

\qffield{Status} Verification: partly checked 2026-10-16 (against the text layer of a local copy of \texttt{alpha-harris}, 2002 Oxford printing, pp.~32--38, 139--144, 195--196 and 400--405). Stability: draft. Currency: concepts timeless; institutional examples as of 2002.
\end{qfdefinition}

\begin{qfdefinition}[Market structures and venues: forums and access, exchanges, OTC and alternative systems, market information systems and transparency, clearing and settlement, and brokered credit]\label{def:market-structures-venues}
\qffield{Statement} ``The trading rules and the trading systems used by a market define its market structure. They determine who can trade; what they can trade; and when, where, and how they can trade. They also determine what information traders can see about orders, quotations, and trades; when they can see it; and who can see it'' \citep[chap.~5, p.~89]{alpha-harris}. Harris classifies a structure by its trading sessions, its execution system and its market information systems (p.~89; sessions and execution systems are \hyperref[def:trading-sessions-execution-systems]{trading sessions and execution systems}), with a minimum price increment (p.~90; \hyperref[def:tick-size]{tick size}). Forums are \emph{physically convened} (pits, posts) or \emph{distributed access} (telephone or screen); some jurisdictions require all trading in an instrument at one exchange (p.~90). \emph{Exchanges} provide forums where members trade; some only host negotiation, others match orders by rules; corporate ownership replaced member ownership at many exchanges around 2000 (chap.~3, sec.~3.2.1, pp.~34--35). In \emph{over-the-counter} markets dealers and brokers arrange trades away from exchanges (p.~35). Harris's 2002 terms for U.S.\ equities: ECNs are order-driven systems not regulated as exchanges; the \emph{third market} is dealers trading listed stocks away from the exchange; alternative trading systems, ECNs among them, trade in the \emph{fourth market} (pp.~35 and~49). Exchanges are search engines: they ``collect information about who wants to trade'' and organise it so that traders find the best opportunities \citep[chap.~19, box, p.~396]{alpha-harris}.

\qffield{Notation} \emph{Pre-trade} (ex ante) and \emph{post-trade} (ex post) transparency: quick reporting of quotes and orders, and of trades; \emph{top of book} against \emph{market by price} \citep[chap.~5, sec.~5.4.4, p.~101]{alpha-harris}. Open-book and closed-book markets (sec.~5.4.7, p.~109).

\qffield{Conditions} \emph{Market information systems} (sec.~5.4, pp.~97--109): collection (electronic systems collect data cheaply; floors need reporters; data were over 20 percent of revenues at many exchanges, p.~97), distribution (members often see more and sooner; real-time and delayed services price-discriminate; U.S.\ securities information processors consolidate quotes and trades, p.~98), order routing, order presentation and order books. \emph{Transparency}: ``Traders are often ambivalent about transparency \dots\ those who know the least about market conditions most favor transparency. Those who know the most oppose transparency'' (p.~101); Lehalle and Laruelle call the same tension a paradox: more shared information should improve price formation, but ``each trader fears `information leakage' '' \citep[Introduction, pp.~19--20]{alpha-laruelle}. \emph{Post-trade infrastructure}: a clearing agent matches buyer and seller records (failures are out-trades or DKs); settlement agents exchange cash and securities, usually net; clearinghouses act ``as buyer for every seller and as seller for every buyer'', require margins and position limits, and are mutualised among clearing members; depositories and custodians hold assets; systems that automate clearing and settlement provide straight-through processing \citep[chap.~3, sec.~3.2.2, pp.~35--38]{alpha-harris}.

\qffield{Practice} \emph{Brokered credit.} Harris's sidebar \emph{Multiplying Credit Checks}: with $10^6$ traders and no brokers, each trader must be ready to check every other, $999{,}999{,}000{,}000$ potential relationships; with 100 brokers of 10{,}000 clients each who guarantee their clients' trades, brokers check clients, clients check brokers and brokers check each other: $2{,}009{,}900$, ``500{,}000 times smaller'' \citep[chap.~7, p.~142]{alpha-harris}. These are ordered counts, $n(n-1)$ and $2n+B(B-1)$. The legacy proposition counts unordered pairs: $n(n-1)/2$ against $n+B(B-1)/2$ ($499{,}999{,}500{,}000$ against $1{,}004{,}950$, the same ratio of about 497{,}500). Either form holds when each client uses exactly one broker that guarantees its trades, clients do not check one another or other brokers, every pair of brokers may trade, and $n$ counts clients only (checked by enumeration on a small network, \nolinkurl{slice35-checks.py}). \emph{Reported volume} depends on architecture: a 100-share trade between public traders reports 100 shares at an order-matching market, 200 with one dealer in between and 300 with two, and the World Federation of Exchanges distinguishes trading-system-view from regulated-environment-view counts (chap.~3, box, p.~48). \emph{RFQ against CLOB.} In fixed income the dealing desk requests quotes from several market makers and takes the best, while the most liquid products trade on electronic books close to equity ones \citep[chap.~2, sec.~2.6, pp.~190--191]{alpha-laruelle}; the protocols are \hyperref[def:fixed-income-market-structure]{fixed-income market structure}. Electronification has not produced one architecture: futures are almost unfragmented, which Lehalle and Laruelle attribute to post-trade complexity, and U.S.\ options are fragmented while European ones are not, because of fragmented clearing (pp.~189--190). \emph{Which venue for which order}: small traders like wide exposure, while large traders want exposure only to traders who will commit to trade with them, through crossing networks, negotiated size or brokers (Harris, chap.~26, sec.~26.3.1, pp.~530--531); the legacy \emph{Architecture Choice} note (small urgent orders to visible venues, large or information-sensitive orders to dark, negotiated or sliced execution) restates this.

\qffield{Pitfalls} Harris's venue names, regulatory classes and settlement cycle (T+3 in 2002, with T+1 planned for 2005, p.~36) are historical; current U.S.\ and EU venue types are in \hyperref[rule:us-equity-market-structure]{U.S.\ equity market structure rules} and \hyperref[rule:eu-trading-venue-transparency]{EU trading venues and transparency}. The legacy taxonomy of market centres (consolidated, fragmented, segmented) is Harris's chap.~26 vocabulary, treated in \hyperref[mech:venue-competition-fragmentation]{venue competition and fragmentation}. Its statement that dark pools, internalisers and crossing networks ``can weaken public price discovery'' is \hyperref[def:dark-pools-crossing]{dark pools and crossing networks}. The legacy designated-market-maker list (single specialist, multiple designated market makers, third-market dealer, pure CLOB) follows Harris's chap.~24 \citep[sec.~24.1, pp.~495--496]{alpha-harris}: specialist systems impose obligations to supply liquidity, usually one specialist per stock; designated multiple-market-maker systems oblige several, obligations being ``best enforced in electronic markets''; third-market dealers' obligations are voluntary.

\qffield{See also} \hyperref[def:trading-sessions-execution-systems]{trading sessions and execution systems}; \hyperref[def:market-intermediaries]{market intermediaries}; \hyperref[def:floor-automated-trading-systems]{floor and automated trading systems}; \hyperref[mech:otc-clearing-collateral]{OTC clearing and collateral}; \hyperref[data:market-data-levels]{market data levels}.

\qffield{Further reading} Harris, chaps.~3, 5, 7 and~24. The legacy file is kept verbatim in \nolinkurl{notes/artefacts/execution-market-making/}; the credit count in \nolinkurl{slice35-checks.py}.

\qffield{Index} market structure; physically convened markets; distributed access; exchanges; demutualisation; over-the-counter; ECN; alternative trading system; third market; fourth market; market information systems; securities information processor; transparency; pre-trade transparency; post-trade transparency; open book; clearing; settlement; net settlement; clearinghouse; central counterparty; custodian; depository; straight-through processing; brokered credit; credit checks; reported volume; RFQ; CLOB; specialists; designated market makers

\qffield{Status} Verification: partly checked 2026-10-16 (against the text layers of local copies of \texttt{alpha-harris}, 2002 Oxford printing, pp.~32--38, 45--49, 89--111, 139--144, 396 and 494--496, and \texttt{alpha-laruelle}, 2018 World Scientific printing, pp.~19--20 and 189--191; the credit counts recomputed). Stability: draft. Currency: concepts timeless; venue classes and settlement cycles as of 2002.
\end{qfdefinition}

\begin{qfdefinition}[Trading sessions and execution systems: continuous and call markets, quote-driven, order-driven, brokered and hybrid markets, auctions and fixings, and specialists]\label{def:trading-sessions-execution-systems}
\qffield{Statement} \emph{Sessions.} ``In continuous trading, traders can attempt to arrange their trades whenever the market is open. In call markets, all trades take place only when the market is called'' \citep[chap.~5, sec.~5.1, p.~89]{alpha-harris}. A call may cover all instruments at once or call them in \emph{rotation}, by 2002 rare; many continuous markets open with a call and restart with one after halts; governments sell debt in calls; some markets (Deutsche Börse, Euronext Paris in 2002) trade their least active stocks only in calls (sec.~5.2.2, pp.~90--91). Calls focus attention ``at the same time and place''; continuous markets serve impatient traders (sec.~5.2.3, p.~91). Some markets run after-hours sessions (sec.~5.2.4, p.~92). \emph{Execution systems}, ``the defining characteristic of a market'' (sec.~5.3, p.~92):
\begin{itemize}
\item \emph{quote-driven} (dealer) markets: in pure form dealers take part in every trade and choose their clients; interdealer brokers let dealers trade anonymously (sec.~5.3.1, pp.~92--93);
\item \emph{order-driven} markets: order precedence rules decide who trades with whom and pricing rules the price; most are auction markets, by single-price auction, continuous two-sided auction or crossing network; auctions are oral (open outcry) or rule-based; traders cannot choose counterparties, so these markets need strong credit arrangements (sec.~5.3.2, pp.~94--95);
\item \emph{brokered} markets: brokers search for \emph{concealed} traders, who will not expose orders, and \emph{latent} traders, who do not know they want to trade until shown an opportunity; typical for blocks, real estate and going concerns (sec.~5.3.3, pp.~95--96);
\item \emph{hybrid} markets mix these, as the 2002 NYSE (order-driven, with specialists obliged to supply liquidity) and Nasdaq (quote-driven, displaying public limit orders) did (sec.~5.3.4, p.~96).
\end{itemize}

\qffield{Conditions} \emph{Auctions and fixings.} A single-price double-sided auction ranks bids and offers and usually maximises feasible volume; securities with little natural interest trade mainly in periodic auctions, also called fixings; opening and closing auctions set reference prices used for margins, fund valuations, cash-settled derivatives and some mergers, which invites attempts to ``mark the close'' \citep[chap.~2, sec.~2.4, pp.~17--18]{s6-hasbrouck}. Design details matter: random stopping times discourage last-instant orders (Euronext), deadlines may be extended after large indicated moves, and large or destabilising orders and the cancellation of stabilising orders may face earlier deadlines (p.~18). Clearing rules and incidence are \hyperref[mech:order-matching-pricing]{order matching and trade pricing}; primary markets' hold on auctions and its effect on venue shares is \hyperref[meth:venue-market-share]{venue market share}. \emph{Specialists} play three roles, as dealers trading for their own account, brokers representing other brokers' orders, and exchange officials responsible for orderly markets, with the auctioneer's role under the third \citep[chap.~24, sec.~24.1, p.~496; secs.~24.2--24.4, pp.~496--501]{alpha-harris}.

\qffield{Practice} \emph{The London gold price as a sequential call.} From 12 September 1919 a daily gold price was set by five founding members, chaired by N~M Rothschild \& Sons, by telephone for the first days and then at a daily meeting \citep[retrieved 2026-10-16]{lbma-gold-centenary}; LBMA describes that arrangement, ``an exclusive membership of five trading houses agreeing the daily price'', as no longer valid, and on 20 March 2015 announced ICE Benchmark Administration as the new independent administrator of the auction (same page). As retrieved on 16 October 2026, IBA runs electronic auctions for gold at 10:30 and 15:00 London time in rounds of fixed duration (for example 30 seconds): it publishes a price, participants enter, change or cancel quantities, and the round ends the auction if the buy--sell imbalance is within a preset threshold; otherwise the price is adjusted and a new round starts; any residual imbalance is shared equally among direct participants \citep[retrieved 2026-10-16]{iba-precious-metals-auctions}. The legacy case study's ``participants submit net interest, revise after tentative prices, and stop when net demand and supply balance'' describes this procedure; it produces one price, not continuous immediacy. \emph{Why calls suit large orders}: a call concentrates opposite interest at one price, and Harris notes that trading halts which reopen with a call switch the pricing rule from discriminatory to uniform, protecting standing limit orders \citep[chap.~28, pp.~572--573]{alpha-harris}.

\qffield{Pitfalls} Session hours, call rotations and the hybrid examples are as of 2002 (Harris) and 2007 (Hasbrouck). The legacy ``auction session'' list (opening, intraday, closing, volatility calls) is a classification; venue rules differ and were not checked. ``More transparency improves price discovery only when participants cannot cheaply game the signal'' is the legacy notes' synthesis; Hasbrouck's deadline rules (p.~18) are the documented countermeasures. ``Floor and brokered systems \dots\ weakens equal access'' merges into \hyperref[def:floor-automated-trading-systems]{floor and automated trading systems}. The gold price's 1919--2015 procedure is from LBMA's history page; the IBA methodology and participant list change over time.

\qffield{See also} \hyperref[def:market-structures-venues]{market structures and venues}; \hyperref[mech:order-matching-pricing]{order matching and trade pricing}; \hyperref[def:floor-automated-trading-systems]{floor and automated trading systems}; \hyperref[mech:circuit-breakers-trading-halts]{circuit breakers and trading halts}; \hyperref[def:dealer-inventory-risk]{dealer inventory risk}.

\qffield{Further reading} Harris, chaps.~5, 6 and~24; Hasbrouck, chap.~2. The legacy file is kept verbatim in \nolinkurl{notes/artefacts/execution-market-making/}.

\qffield{Index} trading sessions; continuous market; call market; rotation; after-hours trading; execution systems; quote-driven market; dealer market; interdealer broker; order-driven market; auction market; open outcry; brokered market; concealed traders; latent traders; hybrid market; fixing; closing auction; marking the close; random end; specialists; London gold fixing; LBMA Gold Price; IBA auction

\qffield{Status} Verification: partly checked 2026-10-16 (against the text layers of local copies of \texttt{alpha-harris}, 2002 Oxford printing, pp.~89--96, 494--501 and 572--573, and \texttt{s6-hasbrouck}, 2007 Oxford printing, pp.~17--19; LBMA's centenary page and IBA's auction page as retrieved on 16 October 2026). Stability: draft. Currency: structures timeless; venue examples as of 2002--2007; gold auction as of October 2026.
\end{qfdefinition}

\begin{qfdefinition}[Liquidity and market quality: liquidity as bilateral search, immediacy, width, depth and resiliency, displayed and latent liquidity, the four observable variables, and liquidity risk]\label{def:liquidity-dimensions}
\qffield{Statement} Harris: ``Liquidity is the ability to trade large size quickly, at low cost, when you want to trade'' \citep[chap.~19, p.~394]{alpha-harris}. It ``is the object of a bilateral search'' in which buyers look for sellers and sellers for buyers, actively or passively (sec.~19.1, p.~395). Searching has inputs (time) and outputs (price and size), whose trade-offs define three dimensions (sec.~19.2, p.~398):
\begin{itemize}
\item \emph{immediacy}: ``how quickly trades of a given size can be arranged at a given cost'';
\item \emph{width}: ``the cost of doing a trade of a given size'', the spread for small trades and including commissions, the cost per unit of liquidity, also called breadth;
\item \emph{depth}: ``the size of a trade that can be arranged at a given cost'';
\end{itemize}
and a fourth, \emph{resiliency}: ``how quickly prices revert to former levels after they change in response to large order flow imbalances initiated by uninformed traders'' (p.~400). Width and depth are duals (p.~399): with a nondecreasing all-in cost $C(q)$ of trading $q$ over a fixed horizon and $D(c)=\max\{q:C(q)\le c\}$, $q\le D(c)$ if and only if $C(q)\le c$, provided the maximum exists. Lehalle and Laruelle's four observable liquidity variables are traded volume, the quantity at the best limits (the \emph{book size}), the bid--ask spread and volatility \citep[chap.~2, sec.~2.2, pp.~143--148]{alpha-laruelle}.

\qffield{Notation} $C(q)$ all-in cost of size $q$; $D(c)$ executable depth within a cost budget $c$; $B$, $c$ and $L$ in the search rule below: expected improvement from one more search, its direct cost, and the expected loss if the current best match disappears.

\qffield{Conditions} \emph{Duality.} The maximum need not exist: with $C(q)=q$ for $q<1$, $C(q)=2$ for $q\ge1$ and $c=1.5$, the feasible sizes are $[0,1)$, the supremum 1 is not feasible, and $q=1\le\sup$ while $C(1)>c$; the equivalence needs $C$ lower semicontinuous (left-continuous at its jumps) and $C(0)\le c$ (checked, \nolinkurl{slice35-checks.py}). \emph{Definitions differ.} Hasbrouck's ``depth, breadth, and resiliency'' uses breadth for a market with many participants, none with market power \citep[chap.~1, sec.~1.2, pp.~4--5]{s6-hasbrouck}, not Harris's cost per unit; precise definitions ``only exist in the contexts of particular models'' (p.~4). \emph{Search.} Continue an active search while the expected benefit of another inquiry exceeds its cost; in a bilateral search ``the cost of continuing to search must include the potential loss of your previous best match'' \citep[chap.~19, secs.~19.1.1--19.1.2, pp.~395--396]{alpha-harris}, hence the legacy rule, search again if $B>c+L$. It is a one-step comparison; with recall and independent offers it coincides with the reservation-value rule (checked on uniform offers with cost 0.02: reservation 0.80, net value 0.8005 against 0.784 and 0.792 for earlier and later stopping), but it is not a general optimal-stopping theorem. Harris's camera example: one more search has expected gain below 2.50 (a gain of at most 10 with probability below 0.25) and costs more than 2.50 at 30 an hour for five minutes, so the buyer stops (box, p.~395).

\qffield{Practice} \emph{Search strategies}: active searchers buy speed with exposure, passive searchers wait and risk not trading, and large traders search actively only among traders who post firm quotes or become passive and undisplayed (Harris, pp.~396--397); the legacy ``selective search'' (trusted counterparties first, then wider) is how brokers use what they know about who wants to trade (pp.~397--398). \emph{Displayed and undisclosed liquidity.} Harris's Table~19-1, after a classification by George Sofianos, lists displayed quotes and orders, liquidity not displayed at a venue, displayed indications, liquidity not displayed at broker-dealers or at the buy side, and latent liquidity, each with its search strategy (``take it'', ``offer to trade'', ``ask'', ``use Liquidnet'') (p.~397). \emph{Book size} averages the quantities at the best bid and ask, summing several limits for small-tick instruments; it is neither an upper nor a lower bound on traded volume, because books replenish (a best bid of 20 throughout and 100 trades of 19 in five minutes trade 1{,}900) and a large order sweeps several limits \citep[chap.~2, sec.~2.2, pp.~143--144]{alpha-laruelle}. \emph{What drives liquidity for an order}: why the trader trades (informed traders find less), what is traded (interest), how well values are known, when (closed markets, suspected news, holidays), what others are doing, displayed imbalances, and who is trading (Harris, sec.~19.2, p.~399; sec.~19.1.2, p.~398). \emph{Resiliency} comes from value traders when they are ``well capitalized, well informed, and willing to trade'' (sec.~19.3.4, p.~403). \emph{Time of day}: volume, volatility, spread and book size share an intraday shape; high volume, volatility and spread after the open, a calm middle of the day, rising volume into the close with the spread at its minimum and the book larger (Royal Dutch Shell, 2011--2016, Lehalle and Laruelle, Figure~2.12, pp.~145--148). \emph{Flow reactivity}, Lehalle and Laruelle's ``correlation between the number of trades and high frequency volatility'', measures whether a venue keeps trading when volatility rises (sec.~2.3.3, pp.~156--158).

\qffield{Pitfalls} Measures see one face each: ``liquidity cannot be reduced to past traded volumes'' (Lehalle and Laruelle, p.~143), and the spread misses depth. The legacy caveat ``book size is a lower bound'' is wrong as stated: hidden, reserve and replenishing orders can make executable size exceed the display, and cancellations before an order arrives can make it smaller (\hyperref[def:order-exposure-decisions]{order exposure decisions}). ``Liquidity risk'' as uncertainty about the cost of turning a desired position into a trade, and ``price continuity'' as the absence of jumps and reversals for ordinary order flow, are the legacy notes' definitions, not Harris's; the legacy remark that forcing continuity against informed flow ``only transfers losses to liquidity suppliers'' is an interpretation consistent with adverse selection (\hyperref[res:adverse-selection-spread-component]{adverse-selection component}) and with Harris's analysis of halts (\hyperref[mech:circuit-breakers-trading-halts]{circuit breakers}). The legacy ``market quality'' list (liquidity, tightness, depth, resiliency, transparency, volatility control and informational efficiency) adds transparency and volatility control to Harris's five characteristics (\hyperref[def:market-functions]{what markets do}); its claim that these criteria conflict is Harris's (chap.~9, p.~218).

\qffield{See also} \hyperref[res:spread-determinants]{spread determinants} (spread and book size when the tick binds; the bounce statistic); \hyperref[meth:spread-volatility-regression]{spread--volatility regression}; \hyperref[def:price-impact]{price impact}; \hyperref[meth:intraday-deseasonalisation]{intraday deseasonalisation}; \hyperref[def:market-intermediaries]{market intermediaries}; \hyperref[def:spread-measures]{spread measures}.

\qffield{Further reading} Harris, chap.~19; Hasbrouck, chap.~1; Lehalle and Laruelle, secs.~2.2--2.3. The legacy file is kept verbatim in \nolinkurl{notes/artefacts/execution-market-making/}; the duality and search checks in \nolinkurl{slice35-checks.py}.

\qffield{Index} liquidity; bilateral search; unilateral search; active search; passive search; stopping rule; reservation value; immediacy; width; breadth; depth; resiliency; depth--width duality; displayed liquidity; undisclosed liquidity; latent liquidity; indications of interest; book size; traded volume; bid--ask spread; volatility; intraday seasonality; flow reactivity; liquidity risk; price continuity; market quality

\qffield{Status} Verification: partly checked 2026-10-16 (against the text layers of local copies of \texttt{alpha-harris}, 2002 Oxford printing, pp.~394--409; \texttt{s6-hasbrouck}, 2007 Oxford printing, pp.~4--6; and \texttt{alpha-laruelle}, 2018 World Scientific printing, pp.~143--149 and 156--158; the duality and the search rule computed). Stability: draft. Currency: concepts timeless; Lehalle and Laruelle's data 2008--2016.
\end{qfdefinition}
